% The volume-uniform Poincare walls — machine-checked obstruction paper.
% Isolated W-3 review lane, reconstructed from the 64969b35 seal.
% REGIME: tricotomy; negative results presented as theorems, not failures;
% NO scores, NO desk language, NO commit chronology beyond the
% reproducibility freeze.  The first artifact remains frozen at source
% checkpoint 12ca1a87 (sealed 64969b35); this is a separate review.
\documentclass[11pt]{article}
\usepackage[margin=1.1in]{geometry}
\usepackage{amsmath,amssymb,amsthm}
\usepackage{booktabs,longtable,array}
\usepackage[colorlinks=true,linkcolor=blue,urlcolor=blue,citecolor=blue]{hyperref}

\newtheorem{theorem}{Theorem}[section]
\newtheorem{lemma}[theorem]{Lemma}
\newtheorem{proposition}[theorem]{Proposition}
\newtheorem{corollary}[theorem]{Corollary}
\newtheorem{remark}[theorem]{Remark}
\theoremstyle{definition}
\newtheorem{definition}[theorem]{Definition}

\emergencystretch=3em
\tolerance=2000

\newcommand{\lean}[1]{\texttt{#1}}
\newcommand{\R}{\mathbb{R}}
\newcommand{\Z}{\mathbb{Z}}
\newcommand{\SU}{\mathrm{SU}}
\newcommand{\su}{\mathfrak{su}}
\newcommand{\CP}{C_{\mathrm{P}}}
\newcommand{\inn}[2]{\langle #1,\, #2\rangle}

\title{The Volume-Uniform Poincar\'e Walls:\\
Machine-Checked Obstructions for Flat and Fluctuation-Sector\\
Block-Poincar\'e Routes\\
to Combes--Thomas Coercivity in Lattice Yang--Mills}
\author{Lluis Eriksson}
\date{July 14, 2026 (revised version)}

\begin{document}
\maketitle

\begin{abstract}
Inside a Lean~4 formalization programme for four-dimensional $\SU(N_c)$
lattice Yang--Mills, we report two machine-checked \emph{negative}
results and the machine-checked infrastructure that makes them
meaningful.  The positive substrate is: (i) a fixed-volume Combes--Thomas chain for
self-adjoint coercive finite-range lattice operators, instantiated on the
flat gauge-fixed covariance of the physical shell, with coercivity
constant $c = \min(1,a)/\CP$ fed by a proved fixed-volume flat
Hodge/block-Poincar\'e inequality; and (ii) the concrete adjoint model of
$\SU(n)$ --- $\su(n)$ with the trace inner product,
$\dim_{\R}\su(n)=n^2-1$, and the isometric transport to Euclidean
coordinates --- so that the abstract adjoint-model interface has a
concrete nontrivial inhabitant and the flat-lane results can be
instantiated with the genuine matricial adjoint model.  The first wall
states that, under the block normalization actually used by the formalized chain, every flat
Hodge/block-Poincar\'e constant obeys $L^d/L^2 \le \CP$ on the fine torus
of side $LN'$, hence the volume-uniform Poincar\'e gate is \emph{provably
false} for $d\ge 3$, $N_c\ge 2$, and no positive coercivity constant
survives all volumes through this route.  The route consumed by the
fixed-volume endpoint is therefore closed, by theorem rather than by
impression.  The original interface left the constant-sector quotient
gate open.  A transverse half-period square mode now supplies a second,
strictly delimited wall.  For $d\ge3$, $N_c\ge2$, arbitrary adjoint model,
and fixed coarse side $N'=2$, the formalization proves
\[
 M\inn{A_M}{K_0A_M}=4\|A_M\|^2,
 \qquad M\|QA_M\|^2\le\|A_M\|^2,
\]
hence $M(\inn{A_M}{K_0A_M}+\|QA_M\|^2)\le5\|A_M\|^2$ and the
volume-uniform quotient gate is false at $N'=2$.
This does not establish the corresponding statement for any other
coarse side.  Everything stated here is checked
by Lean~4 against a pinned Mathlib, with zero \lean{sorry}, zero project
axioms, and a committed axiom-oracle transcript; a theorem--artifact map
and reproduction instructions are included.  The results close two
routes under the current flat normalization; they establish neither a
necessary condition for Yang--Mills nor any part of the continuum
construction or mass-gap theorem.
\end{abstract}

\section{Introduction}

Constructive approaches to the four-dimensional Yang--Mills mass gap in
the Balaban tradition~\cite{Balaban84a,Balaban84b} route spectral
information through renormalization-group covariances whose kernels must
decay \emph{uniformly in the volume}.  A standard engine for such decay
is the Combes--Thomas argument~\cite{CombesThomas}: a self-adjoint
coercive finite-range operator has an inverse whose kernel entries decay
exponentially in operator norm, at a rate governed by the coercivity
constant.  Whoever supplies
the coercivity therefore controls the physics; and in a gauge-fixed
lattice formalization the natural supplier is a Poincar\'e-type
inequality tying the field norm to a Hodge term plus a block-averaging
constraint.

This paper reports what happened when that plan was executed
\emph{formally}, in Lean~4~\cite{Lean4} against
Mathlib~\cite{mathlib}, inside a long-running formalization programme
whose earlier machine-checked artifacts include a volume-uniform
Wilson-loop area law via a formalized cluster
expansion~\cite{ErikssonVU}, rooted-tree majorants for polymer expansions
with holes~\cite{ErikssonHoles}, a mechanized Bessel-ratio calculus for
lattice gauge expansions~\cite{ErikssonC4,ErikssonC3,ErikssonC2}, an
exact two-dimensional $\SU(2)$ theory~\cite{Eriksson2D}, and the audited
programme map~\cite{ErikssonMap}.

The outcome has four layers, and we state them with the labels they
deserve.

\textbf{Positive (fixed volume).}  A complete Combes--Thomas chain
for coercive finite-range operators on finite metric graphs,
instantiated at the flat gauge-fixed covariance of the physical shell:
for every fixed volume there exist a positive rate $\theta$ and amplitude
$2/c$ with
$\bigl\|(K^{-1})_{q,p}\bigr\|_{\mathrm{op}} \le (2/c)\,e^{-\theta\, d(q,p)}$
for all bond pairs $(q,p)$,
where $c = \min(1,a)/\CP$ and $\CP$ is a \emph{proved} fixed-volume flat
Hodge/block-Poincar\'e constant (Section~\ref{sec:ct}).

\textbf{Positive (non-vacuity).}  The abstract adjoint-model interface
consumed by the chain is inhabited by the \emph{true} matricial adjoint
model of $\SU(n)$: $\su(n)$ as the real inner-product space of traceless
skew-Hermitian matrices under the trace form, its real dimension
$n^2-1$, and the isometric transport to
$\R^{n^2-1}$ carrying $\mathrm{Ad}(g)X = gXg^{\dagger}$
(Section~\ref{sec:adj}).  The flat-lane results also admit the trivial
inhabitant, since only $\mathrm{Ad}(1)$ is consumed there; the matricial
construction removes the intended-model vacuity and prepares the
nontrivial-background lane.

\textbf{Negative (the constant-sector wall).}  Under the block
normalization actually used by the formalized chain --- the unscaled
line-integral block map, which sends constants to $L$ times constants ---
\emph{every} flat Hodge/block-Poincar\'e constant obeys
\[
\frac{L^d}{L^2} \;\le\; \CP
\]
on the fine torus of side $N = LN'$.  Consequently (Section~\ref{sec:wall}):
the volume-uniform Poincar\'e gate is \emph{false} for $d \ge 3$,
$N_c\ge 2$; the Combes--Thomas coercivity fed by this route obeys
$c \le \min(1,a)\,L^2/L^d \to 0$; and no single $c_0>0$ is dominated at
every volume by the route's coercivity.  The per-volume constants
\emph{do} exist, so the obstruction kills a real route, not an empty one.

\textbf{Negative (the fluctuation-sector wall at $N'=2$).}
The constant-sector quotient interface (Section~\ref{sec:quotient})
excludes the first witness and, at the interface stage, left the
volume-uniform quotient gate open.  Section~\ref{sec:falsifier} evaluates
a different witness: a transverse square-wave fluctuation mode on fine
side $2M$, paired with the block map from side $2M$ to the fixed coarse
side $2$.  Its Hodge contribution is exactly $4/M$ times its norm
squared and its block contribution is at most $1/M$ times its norm
squared for $d\ge3$.  The resulting $5/M$ bound contradicts a constant
chosen before the scale quantifier.  Thus
$\lean{VolumeUniformQuotientPoincareGate}\ d\ 2\ N_c\ \rho$ is false for
$d\ge3$ and $N_c\ge2$.  No statement for another coarse side is inferred.

\subsection*{Honest scope}

This paper proves no mass gap, no volume-uniform Combes--Thomas bound,
and nothing about the thermodynamic or continuum limits.  Neither gate
is shown to be a necessary, equivalent, or exhaustive condition for the
Yang--Mills existence-and-gap problem~\cite{JaffeWitten}.  In particular,
the second wall is restricted literally to $N'=2$, $d\ge3$, $N_c\ge2$,
arbitrary adjoint model, and the current flat unscaled block operator
with its present coarse norm.  It gives no theorem for all coarse sizes,
for a rescaled or weighted block operator, or for the interacting Wilson
Hessian.  The results are navigational obstructions: they close two
specific flat routes.  They are not a quantitative measure of progress
toward the Clay problem, and no percentage-of-solution interpretation is
attached to them.

\section{The formal setting}\label{sec:setting}

All objects below are Lean definitions in the repository accompanying
this paper (Section~\ref{sec:repro}); we give their mathematical content
and their Lean names.

\paragraph{Cochains.}  Fix $d,N,N_c\in\mathbb{N}$, all positive.  Sites
form the discrete torus $\Lambda_N=(\Z/N)^d$
(\lean{FinBox d N}); positive bonds are pairs $b=(x,k)$ of a site and a
direction (\lean{PhysicalBond d N}).  A \emph{physical one-cochain} is a
map $A$ from positive bonds to the internal coordinate space
$\R^{N_c^2-1}$ (\lean{SUNLieCoord Nc}), with the $\ell^2$ inner product
in both the bond and internal indices
(\lean{PhysicalGaugeOneCochain d N Nc}).

\paragraph{The block map.}  For $N=LN'$, the flat block constraint
$Q\colon$ fine cochains (side $LN'$) $\to$ coarse cochains (side $N'$)
assigns to a coarse bond the \emph{unscaled} line integral (sum) of the
$L$ fine bond values along the corresponding block line
(\lean{flatBlockConstraintQCLM}, built on \lean{linAvg}).  The
normalization fact that drives everything below is exact
(\lean{flatBlockConstraintQCLM\_constant}):
\begin{equation}\label{eq:Qconst}
Q\,(A_v) \;=\; L\cdot A^{\mathrm{coarse}}_v,
\end{equation}
where $A_v$ denotes the direction-wise constant cochain
$A_v(x,k) = v_k$ for a fixed family $v\colon \{1,\dots,d\}\to\R^{N_c^2-1}$
(\lean{constantPhysicalGaugeOneCochain}).

\paragraph{The flat Hodge operator.}  At the trivial gauge background,
$K_0 = D_1^{\dagger}D_1 + D_0 D_0^{\dagger}$
(\lean{flatGaugeHodgeK0CLM}), so that
$\inn{A}{K_0 A} = \|\mathrm{curl}\,A\|^2 + \|\mathrm{div}\,A\|^2$ with
the standard lattice curl on plaquettes and backward divergence on sites
(\lean{flatGaugeHodgeK0\_inner\_right}).  Constants are flat-harmonic:
$K_0 A_v = 0$
(\lean{flatGaugeHodgeK0CLM\_constantPhysicalGaugeOneCochain}).

\paragraph{The Poincar\'e predicate.}  For an adjoint model $\rho$
(Section~\ref{sec:adj}) and $\CP\in\R$,
\lean{FlatGaugeHodgePoincare d L N' Nc $\rho$ CP} is the statement
\begin{equation}\label{eq:poincare}
0 < \CP \quad\wedge\quad
\forall A:\ \|A\|^2 \;\le\; \CP\,\bigl(\inn{A}{K_0A} + \|QA\|^2\bigr),
\end{equation}
over fine cochains at side $LN'$.  The repository proves, at every fixed
volume, that such constants exist
(\lean{exists\_flatGaugeHodgePoincare}); this non-vacuity is what makes
the negative results of Section~\ref{sec:wall} statements about a live
route.

\section{The fixed-volume Combes--Thomas endpoint}\label{sec:ct}

The Combes--Thomas chain is formalized for abstract self-adjoint
coercive finite-range operators on finite metric graphs: the tilt
algebra and tilted entry bounds, the block Schur bound, coercivity
survival under tilting at the budget
$M(e^{\theta R}-1)N_R \le c/2$, the tilted inverse with
$\|K_\theta^{-1}\|\le 2/c$, and the kernel extraction at a root
(\lean{PhysicalCoerciveCombesThomas*.lean}).  Instantiated at the flat
gauge-fixed covariance of the physical shell (stencils and locality
proved, not assumed), the endpoint reads:

\begin{theorem}[fixed-volume Combes--Thomas;
\lean{flatGaugeFixedCovariance\_CT\_fixedVolume},
\lean{exists\_flatGaugeFixedCovariance\_CT\_fixedVolume}]
Fix the volume parameters $(d,L,N')$, a mass $a>0$, and a flat
Hodge/block-Poincar\'e constant $\CP$ as in \eqref{eq:poincare}.  Let
$c=\min(1,a)/\CP$.  Then there exists $\theta>0$ meeting the explicit
tilt budget such that for all bond pairs $(q,p)$ and every internal
vector $v$,
\[
\bigl\|(K^{-1})_{q,p}\, v\bigr\| \;\le\; \frac{2}{c}\;
e^{-\theta\, d(q,p)}\,\|v\|,
\]
equivalently
$\|(K^{-1})_{q,p}\|_{\mathrm{op}} \le (2/c)\,e^{-\theta\,d(q,p)}$ ---
the form quantified literally by
\lean{PhysicalCovarianceExponentialKernelBound}.
\end{theorem}

Fixed volume means fixed volume: $\CP$, hence $c$ and $\theta$, may
depend on $(d,L,N',a)$.  The question the rest of this paper answers is
whether \emph{this} $c$ can be made volume-uniform.  It cannot.

\section{The true adjoint model of $\SU(n)$}\label{sec:adj}

The chain consumes an abstract interface \lean{SUNAdjointModel n}: a
family of isometries $\mathrm{Ad}(g)$ of $\R^{n^2-1}$ indexed by
$g\in\SU(n)$ with $\mathrm{Ad}(1)=\mathrm{id}$,
$\mathrm{Ad}(gh)=\mathrm{Ad}(g)\mathrm{Ad}(h)$, and inner-product
invariance.  A trivial inhabitant ($\mathrm{Ad}(g)=\mathrm{id}$) exists
and suffices for the flat lane, where only $\mathrm{Ad}(1)$ is consumed;
leaving matters there would invite a vacuity objection the moment a
non-trivial background appears.  The repository therefore constructs the
genuine object:

\begin{theorem}[the matricial adjoint model;
\lean{matrixSUNAdjointModel}]\label{thm:adj}
Let $\su(n)$ be the real vector space of traceless skew-Hermitian
$n\times n$ complex matrices with the trace form
$\inn{X}{Y} = \operatorname{Re}\operatorname{tr}(X^{\dagger}Y)$
(\lean{SUNAdjointMatrixSubstrate}, \lean{SUNAdjointInnerSpace}).  Then
$\dim_{\R}\su(n) = n^2-1$ (\lean{finrank\_suLie}), the reindexed
orthonormal-basis representation gives a linear isometric equivalence
$\su(n)\simeq\R^{n^2-1}$ (\lean{suLieCoordIso}), and conjugating the
adjoint action $X\mapsto gXg^{\dagger}$ by this isometry yields a
\lean{SUNAdjointModel n} whose action laws and inner-product invariance
are proved, not hypothesized.
\end{theorem}

The dimension count is itself formalized (Hermitian/skew decomposition,
multiplication-by-$i$ equivalence, ambient real dimension $2n^2$, and
rank--nullity through the imaginary-trace functional with explicit
surjectivity witness $i\cdot\mathbf{1}$).  The coordinate choice is the
non-canonical standard orthonormal basis; any other choice conjugates
the model by a fixed orthogonal map, and nothing downstream depends on
it.

\section{The wall}\label{sec:wall}

Everything in this section is proved in
\lean{PhysicalGaugeFlatPoincare.lean} (the fixed-volume audit) and
\lean{PhysicalPoincareWall.lean} (the volume-uniform obstruction).

\subsection{The constant sector forces the constant}

\begin{lemma}[exact sector normalization;
\lean{flatBlockConstraint\_controls\_constantSector}]\label{lem:sector}
For every $v$,
\[
\|A_v\|^2 \;=\; \frac{L^d}{L^2}\,\|Q A_v\|^2 .
\]
\end{lemma}

\begin{proof}
$\|A_v\|^2 = N^d\sum_k\|v_k\|^2$ with $N=LN'$
(\lean{norm\_sq\_constantPhysicalGaugeOneCochain}), while by
\eqref{eq:Qconst},
$\|QA_v\|^2 = L^2 N'^d \sum_k\|v_k\|^2$
(\lean{flatBlockConstraintQCLM\_constant\_norm\_sq}); divide.
\end{proof}

\begin{theorem}[the wall, fixed volume;
\lean{flatPoincare\_constantSector\_lower\_bound}]\label{thm:wall}
Let $N_c\ge 2$ and let $\CP$ satisfy \eqref{eq:poincare}.  Then
\[
\frac{L^d}{L^2}\;\le\;\CP .
\]
\end{theorem}

\begin{proof}
For $N_c \ge 2$ the internal space $\R^{N_c^2-1}$ is nontrivial, so
there is $v$ with $\sum_k \|v_k\|^2>0$
(\lean{exists\_pos\_constantSector}; for $N_c=1$, $\su(1)=0$ and the
sector is empty --- the hypothesis is stated, not hidden).  Apply
\eqref{eq:poincare} to $A_v$: the Hodge term vanishes ($K_0A_v=0$), so
$\|A_v\|^2 \le \CP\|QA_v\|^2$, and Lemma~\ref{lem:sector} gives the
claim after cancelling $\|QA_v\|^2>0$.
\end{proof}

\begin{corollary}[linear form;
\lean{flatPoincare\_linear\_lower\_bound}]\label{cor:linear}
If moreover $d\ge 3$, then $L \le \CP$.  (For $d=4$ the forced growth is
$L^2$; the linear form is all the divergence the obstruction needs.
$d=2$ is the scale-invariant exemption $L^d/L^2=1$ and is deliberately
not claimed.)
\end{corollary}

\subsection{The route's coercivity dies with the volume}

\begin{theorem}[coercivity upper bound;
\lean{flatPoincare\_coercivity\_upper\_bound}]
For $a>0$, $N_c\ge2$, and any $\CP$ as in \eqref{eq:poincare},
\[
\frac{\min(1,a)}{\CP} \;\le\; \min(1,a)\,\frac{L^2}{L^d}.
\]
\end{theorem}

\begin{theorem}[no volume-uniform coercivity through flat Poincar\'e]
\label{thm:noc}
{\raggedright\emph{Lean:}
\path{no_volumeUniform_coercivity_via_flatPoincare}.\par}
Let $d\ge3$, $N_c\ge2$, $a>0$.  Suppose $c_0\in\R$ is dominated at every
volume: for every $L\in\mathbb{N}$ there exists $\CP(L)$ satisfying
\eqref{eq:poincare} at side $(L+1)N'$ with
$c_0 \le \min(1,a)/\CP(L)$.  Then $c_0\le 0$.
\end{theorem}

\begin{proof}
If $c_0>0$, choose (Archimedes) $n$ with $n > \min(1,a)/c_0$.  By
Corollary~\ref{cor:linear}, $\CP(n)\ge n+1 > \min(1,a)/c_0$, whence
$\min(1,a)/\CP(n) < c_0$, contradicting domination at volume $n+1$.
\end{proof}

\begin{theorem}[the gate is false;
\lean{volumeUniformFlatPoincareGate\_false}]\label{thm:gatefalse}
For $d\ge3$, $N_c\ge2$, there is no single $\CP>0$ satisfying
\eqref{eq:poincare} at every volume $L+1$:
the proposition \lean{VolumeUniformFlatPoincareGate} --- with the
constant quantified \emph{before} the volume --- is provably false.
\end{theorem}

\begin{remark}[non-vacuity of the killed route;
\lean{perVolume\_flatPoincare\_family}]
The hypothesis family of Theorem~\ref{thm:noc} is inhabited: per-volume
constants exist at every volume.  The obstruction closes a route that is
genuinely open at each fixed volume.
\end{remark}

\begin{remark}[quantifier discipline]
The gate places its constant before the volume quantifier.  The reversed
order --- for every volume there exists a constant --- is exactly the
(true) per-volume statement, and conflating the two is the standard way
such obstructions get lost in prose.  The formalization makes the
distinction type-level.
\end{remark}

\section{The constant-sector quotient interface}\label{sec:quotient}

The wall witness lives in one sector.  The natural continuation
restricts the Poincar\'e question to its orthogonal complement --- the
physical fluctuation space.  The repository formalizes the
\emph{interface} (\lean{PhysicalPoincareSectorQuotient.lean}); we use
the word ``interface'' deliberately: no quotient type or packaged
orthogonal-complement submodule is constructed, and the predicate is a
restriction to fluctuation cochains.

\begin{definition}[fluctuation cochains; \lean{IsFluctuationCochain}]
$A$ is a fluctuation cochain if $\inn{A_v}{A}=0$ for every constant
generator $A_v$ --- equivalently, $A$ is orthogonal to the span of the
constant sector.  (The generator-wise form is chosen for formal
robustness.)
\end{definition}

\begin{proposition}[interface facts]
\emph{(i)} The constant sector is flat-harmonic
(\lean{constantSectorLin\_harmonic}) --- which is why only the block
term sees it, with the normalization of Lemma~\ref{lem:sector}.
\emph{(ii)} The full predicate \eqref{eq:poincare} implies its
restriction to fluctuation cochains at the same constant
(\lean{quotientFlatPoincare\_of\_flatPoincare}); the restricted
statement is strictly weaker, so Theorem~\ref{thm:gatefalse} does
\emph{not} refute the quotient gate.
\emph{(iii)} Restricted constants exist at every fixed volume
(\lean{exists\_quotientFlatPoincare}).
\emph{(iv) (Non-transfer.)}  A nonzero constant cochain is not a
fluctuation cochain (\lean{constant\_not\_fluctuation}): the wall
\emph{witness} is excluded from the fluctuation space by construction.
\end{proposition}

\begin{remark}[what the interface stage established]
At the interface stage the volume-uniform quotient gate
(\lean{VolumeUniformQuotientPoincareGate}) was neither proved nor
refuted.  Non-transfer removed the W-1 witness, not the possibility of a
second wall.  The present revision preserves that historical statement
and adds a different fluctuation witness.  Its endpoint is conclusive
only at fixed coarse side $N'=2$.
\end{remark}

\section{The low-mode falsifier and the second wall at $N'=2$}
\label{sec:falsifier}

A \emph{one-sided falsifier} for the quotient gate needs an explicit
fluctuation family on which the quadratic form of \eqref{eq:poincare}
can be evaluated exactly.  Four modules now implement the complete
chain: construction, Hodge evaluation, block response, and endpoint.

\begin{definition}[the half-period (two-interface) square-wave mode;
\lean{squareModeCochain}]
On even side lengths $N=2M$ ($M\ge1$), with oscillation direction $j$,
bond direction $i$, and internal vector $w$:
\[
A(x,k) \;=\;
\begin{cases}
\varepsilon(x_j)\, w, & k=i,\\[2pt]
0, & k\neq i,
\end{cases}
\qquad
\varepsilon(m)=\begin{cases}+1,& m<M,\\ -1,& m\ge M,\end{cases}
\]
(\lean{squareSign}).  The square profile is chosen over the cosine so
that every identity below is an exact finite $\pm1$ computation; odd
side lengths are out of scope by declared parameterization (a falsifier
needs only \emph{some} volume sequence).
\end{definition}

\begin{proposition}[exact identities]
\emph{(i)} $\sum_m \varepsilon(m) = 0$ (\lean{sum\_squareSign}), hence
$A$ is a fluctuation cochain
(\lean{squareModeCochain\_isFluctuation}) --- exact orthogonality, no
approximation.
\emph{(ii)} $\|A\|^2 = (2M)^d\,\|w\|^2$ exactly
(\lean{norm\_sq\_squareModeCochain}).
\emph{(iii)} For $w\neq0$ the mode equals no constant cochain
(\lean{squareModeCochain\_ne\_constant}), and the fluctuation space is
nontrivial for $N_c\ge2$ (\lean{exists\_nonzero\_fluctuation}).
\end{proposition}

For the endpoint choose distinct directions $i\ne j$, let
$A_M=\lean{squareModeFine}\ d\ M\ N_c\ i\ j\ w$, and apply the block map
with block length $M$ and coarse side $2$.  The fine side is therefore
$M\cdot2=2M$.  The Hodge calculation gives
\begin{equation}\label{eq:w3hodge}
 M\inn{A_M}{K_0A_M}=4\|A_M\|^2.
\end{equation}
The block and fine norms are exact:
\[
 \|QA_M\|^2=2^dM^2\|w\|^2,
 \qquad
 \|A_M\|^2=(2M)^d\|w\|^2.
\]
Thus, when $d\ge3$ and $M\ge1$,
\begin{equation}\label{eq:w3block}
 M\|QA_M\|^2\le\|A_M\|^2.
\end{equation}
Combining \eqref{eq:w3hodge} and \eqref{eq:w3block} without division
produces
\begin{equation}\label{eq:w3total}
 M\bigl(\inn{A_M}{K_0A_M}+\|QA_M\|^2\bigr)
 \le5\|A_M\|^2.
\end{equation}

\subsection*{Dependency record for the four capstone theorems}

\begingroup
\small
\begin{longtable}{@{}>{\raggedright\arraybackslash}p{0.30\textwidth}>{\raggedright\arraybackslash}p{0.62\textwidth}@{}}
\toprule
Lean theorem & Mathematical role and immediate dependencies \\
\midrule
\endfirsthead
\toprule
Lean theorem & Mathematical role and immediate dependencies \\
\midrule
\endhead
\path{squareMode_mul_hodge_eq_four_norm_sq}
& \textbf{Statement:} \eqref{eq:w3hodge}.  \textbf{Inputs:}
\path{flatGaugeHodgeK0_squareMode_energy_mul_side} and the definitional
identification of \path{squareModeFine} with the square-mode cochain.
\textbf{Role:} exact Hodge contribution. \\
\addlinespace
\path{squareMode_mul_block_le_norm_sq}
& \textbf{Statement:} \eqref{eq:w3block}.  \textbf{Hypotheses:}
$d\ge3$, $M\ge1$, $i\ne j$.  \textbf{Inputs:}
\path{norm_sq_flatBlockConstraintQCLM_squareModeFine_of_ne} and
\path{norm_sq_squareModeFine}.  \textbf{Role:} transverse block control. \\
\addlinespace
\path{squareMode_mul_totalEnergy_le_five_norm_sq}
& \textbf{Statement:} \eqref{eq:w3total}.  \textbf{Inputs:} the preceding
two theorems.  \textbf{Role:} division-free Rayleigh-numerator bound. \\
\addlinespace
\texttt{volumeUniformQuotient}\newline
\texttt{PoincareGate\_two\_false}
& \textbf{Statement:}
the quotient gate is false at coarse side $2$.
\textbf{Hypotheses:} $d\ge3$, $N_c\ge2$; $\rho$ arbitrary.
\textbf{Inputs:} \eqref{eq:w3total}, fluctuation membership, exact norm
positivity, a nonzero internal vector, and the Archimedean choice of $M$.
\textbf{Role:} quantifier-respecting endpoint. \\
\bottomrule
\end{longtable}
\endgroup

\begin{theorem}[the quotient gate is false at coarse side two]
\label{thm:quotgatefalse}
{\raggedright\emph{Lean:}
\path{volumeUniformQuotientPoincareGate_two_false}.\par}
For $d\ge3$, $N_c\ge2$, and every adjoint model $\rho$,
\[
 \neg\,\mathrm{Gate}_{\mathrm{quot}}(d,2,N_c,\rho).
\]
\end{theorem}

\begin{proof}
Assume a gate witness $\CP>0$ whose scale quantifier comes after the
choice of $\CP$.  Choose $M\in\mathbb N$ with $M>5\CP$, choose distinct
directions and a nonzero internal vector, and instantiate the gate at the
corresponding fine side $2M$.  The mode is a fluctuation and has positive
norm.  Multiplying the Poincar\'e inequality by $M$ and using
\eqref{eq:w3total} gives
\[
 M\|A_M\|^2\le5\CP\|A_M\|^2.
\]
Positivity permits cancellation, so $M\le5\CP$, a contradiction.  The
formal proof performs this argument without dividing by $\|A_M\|^2$.
\end{proof}

\begin{remark}[scope and one-sided discipline]
This is a \emph{second formally certified wall for the current flat
quotient-Poincar\'e route, at fixed coarse side $N'=2$}.  It does not
establish falsity for another coarse side, necessity or exhaustiveness of
the gate for Yang--Mills, or any statement about the mass gap.  The
falsifier has fired in its conclusive direction; had its quotient stayed
bounded, that fact alone would not have proved the gate.
\end{remark}

\section{Theorem--artifact map}\label{sec:map}

\begingroup
\footnotesize
\begin{longtable}{@{}>{\raggedright\arraybackslash}p{0.30\textwidth}>{\raggedright\arraybackslash}p{0.62\textwidth}@{}}
\toprule
Statement & Lean artifact (module under \lean{YangMills/RG/}) \\
\midrule
\endfirsthead
\toprule
Statement & Lean artifact (module under \lean{YangMills/RG/}) \\
\midrule
\endhead
CT endpoint, fixed volume & \path{flatGaugeFixedCovariance_CT_fixedVolume}\newline
(\lean{PhysicalShellCombesThomasEndpoint}) \\
CT positive-rate witness & \path{exists_flatGaugeFixedCovariance_CT_fixedVolume}\newline
(\lean{PhysicalShellCombesThomasEndpoint}) \\
$\dim_\R \su(n)=n^2-1$ & \path{finrank_suLie} (\lean{SUNAdjointDimension}) \\
Isometric transport & \path{suLieCoordIso} (\lean{SUNAdjointModelInstance}) \\
True adjoint model & \path{matrixSUNAdjointModel} (\lean{SUNAdjointModelInstance}) \\
Sector witness ($N_c\ge2$) & \path{exists_pos_constantSector} (\lean{PhysicalPoincareWall}) \\
Wall, fixed volume & \path{flatPoincare_constantSector_lower_bound} (\lean{PhysicalPoincareWall}) \\
Wall, linear form ($d\ge3$) & \path{flatPoincare_linear_lower_bound} (\lean{PhysicalPoincareWall}) \\
Coercivity upper bound & \path{flatPoincare_coercivity_upper_bound} (\lean{PhysicalPoincareWall}) \\
No volume-uniform $c_0$ & \path{no_volumeUniform_coercivity_via_flatPoincare} (\lean{PhysicalPoincareWall}) \\
Full gate false & \path{volumeUniformFlatPoincareGate_false} (\lean{PhysicalPoincareWall}) \\
Per-volume non-vacuity & \path{perVolume_flatPoincare_family} (\lean{PhysicalPoincareWall}) \\
Fluctuation predicate & \path{IsFluctuationCochain} (\lean{PhysicalPoincareSectorQuotient}) \\
Full $\Rightarrow$ quotient & \path{quotientFlatPoincare_of_flatPoincare} (\lean{PhysicalPoincareSectorQuotient}) \\
Non-transfer of first witness & \path{constant_not_fluctuation} (\lean{PhysicalPoincareSectorQuotient}) \\
Quotient-gate predicate & \path{VolumeUniformQuotientPoincareGate} (\lean{PhysicalPoincareSectorQuotient}) \\
Square profile, zero sum & \path{sum_squareSign} (\lean{PhysicalPoincareLowModeFalsifier}) \\
Mode is fluctuation & \path{squareModeCochain_isFluctuation} (\lean{PhysicalPoincareLowModeFalsifier}) \\
Mode norm, exact & \path{norm_sq_squareModeCochain} (\lean{PhysicalPoincareLowModeFalsifier}) \\
Hodge identity & \path{squareMode_mul_hodge_eq_four_norm_sq} (\lean{PhysicalPoincareLowModeEndpoint}) \\
Block inequality & \path{squareMode_mul_block_le_norm_sq} (\lean{PhysicalPoincareLowModeEndpoint}) \\
Total-energy inequality & \path{squareMode_mul_totalEnergy_le_five_norm_sq} (\lean{PhysicalPoincareLowModeEndpoint}) \\
Quotient gate false at $N'=2$ & \path{volumeUniformQuotientPoincareGate_two_false} (\lean{PhysicalPoincareLowModeEndpoint}) \\
\bottomrule
\end{longtable}
\endgroup

Every listed name is exercised by the axiom-oracle script.  In particular,
all sixteen W-3b/W-3c/endpoint queries, including the four capstone rows,
print exactly \lean{[propext, Classical.choice, Quot.sound]}.

\section{Reproducibility}\label{sec:repro}

\paragraph{Repository.}
\url{https://github.com/lluiseriksson/THE-ERIKSSON-PROGRAMME}
(review branch \lean{poincare-wall-w3-review}).  The first artifact remains
frozen at source checkpoint \lean{12ca1a87}, sealed by
\lean{64969b35}.  The W-3 review is reconstructed from that seal.  Its
isolated Lean and oracle checkpoint is \lean{90ed3585}; no result or
document from a different paper lane is included in the review delta.
The scientific manuscript checkpoint is \lean{009da5b0}.

\paragraph{Toolchain.}  Lean \lean{leanprover/lean4:v4.29.0-rc6};
Mathlib pinned to commit
\lean{0764272048}\ldots{} (lakefile and manifest agree).  Build:
\lean{lake build YangMillsCore} --- expected
\lean{Build completed successfully (8413 jobs)}.  The increase of three
jobs is exact: the falsifier module already belonged to the frozen tree,
while the Hodge, block, and endpoint modules are new.

\paragraph{Axiom oracle.}  \lean{lake env lean oracle\_check.lean}
executes $2210$ \lean{\#print axioms} invocations: the frozen $2194$ plus
sixteen individually identified W-3 queries.  A separate focal run of
those sixteen returns exactly
\lean{[propext, Classical.choice, Quot.sound]} for every target, with
zero \lean{sorryAx} and no project axiom.  The full transcript is
\path{docs/oracle-transcripts/ORACLE-20260714-90ed3585.txt}: $2124$
results use the full trio, $49$ use \lean{[propext, Quot.sound]}, $17$
use \lean{[propext]}, and $20$ are axiom-free.  The LF-normalized raw
output has SHA-256
\path{09044831B051646510F828C6E1C3A45B9BA3503482DB65866C695259E43DD75C}.

\section*{Continuation from this submission}

The Hodge evaluation, block estimate, division-free total-energy bound,
and endpoint at $N'=2$ are now theorems, not future work.  The surviving
questions are deliberately separate: other coarse sides; rescaled or
weighted block operators with their normalization re-audited from first
principles; and coercivity from the interacting Hessian.  No claim that
any of these routes succeeds is made.

\begin{thebibliography}{99}

\bibitem{Balaban84a}
T.~Ba{\l}aban,
\emph{Propagators and renormalization transformations for lattice gauge
theories.~I},
Comm.\ Math.\ Phys.\ \textbf{95} (1984) 17--40.

\bibitem{Balaban84b}
T.~Ba{\l}aban,
\emph{Propagators and renormalization transformations for lattice gauge
theories.~II},
Comm.\ Math.\ Phys.\ \textbf{96} (1984) 223--250.

\bibitem{CombesThomas}
J.~M.~Combes and L.~Thomas,
\emph{Asymptotic behaviour of eigenfunctions for multiparticle
Schr\"odinger operators},
Comm.\ Math.\ Phys.\ \textbf{34} (1973) 251--270.

\bibitem{KoteckyPreiss}
R.~Koteck\'y and D.~Preiss,
\emph{Cluster expansion for abstract polymer models},
Comm.\ Math.\ Phys.\ \textbf{103} (1986) 491--498.

\bibitem{JaffeWitten}
A.~Jaffe and E.~Witten,
\emph{Quantum Yang--Mills theory},
Clay Mathematics Institute Millennium Prize problem description (2000).

\bibitem{Lean4}
L.~de~Moura and S.~Ullrich,
\emph{The Lean 4 theorem prover and programming language},
CADE-28, LNCS 12699 (2021) 625--635.

\bibitem{mathlib}
The mathlib community,
\emph{The Lean mathematical library},
CPP 2020, 367--381.

\bibitem{ErikssonVU}
L.~Eriksson,
\emph{A machine-checked volume-uniform Wilson-loop area law via a
formalized cluster expansion},
ai.viXra.org:2607.0005 (2026).

\bibitem{ErikssonHoles}
L.~Eriksson,
\emph{Machine-checked rooted-tree majorants for polymer expansions with
holes},
ai.viXra.org:2607.0025 (2026).

\bibitem{ErikssonC4}
L.~Eriksson,
\emph{A machine-checked proof of an Amos-type bound for modified Bessel
ratios at real order},
ai.viXra.org:2607.0030 (2026).

\bibitem{ErikssonC3}
L.~Eriksson,
\emph{A machine-checked proof of the Amos bound for modified Bessel
function ratios},
ai.viXra.org:2607.0032 (2026).

\bibitem{ErikssonC2}
L.~Eriksson,
\emph{One Amos bound, three consumer sites: machine-checked
Bessel-ratio calculus for lattice gauge expansions},
ai.viXra.org:2607.0033 (2026).

\bibitem{Eriksson2D}
L.~Eriksson,
\emph{Exact two-dimensional $\SU(2)$ Yang--Mills in Lean: Weyl
integration, heat-kernel convolution, Migdal invariance, and the exact
simple-loop area law},
ai.viXra.org:2607.0035 (2026).

\bibitem{ErikssonCatalan}
L.~Eriksson,
\emph{A machine-verified bijective proof of the rooted child-factorial
Catalan identity over spanning trees of the complete graph},
ai.viXra.org:2607.0001 (2026).

\bibitem{ErikssonMap}
L.~Eriksson,
\emph{The master map --- audit experiments report: mechanical audit
experiments and reproducibility appendix for the 2602-series programme
on 4D $\SU(N)$ Yang--Mills},
ai.viXra.org:2602.0117 (2026).

\end{thebibliography}

\end{document}
